\chapter{OS Security}

Operating system security aims to protect system resources while ensuring correct and controlled access. The fundamental goals remain the same: confidentiality, integrity, and authentication. Since many system resources are sensitive, mechanisms are required to prevent unauthorized access and misuse.

\section{Access Control}

\subsection{Fundamentals}
To enforce security goals, we first need a way to regulate how resources are accessed. Access control is the selective restriction of access to resources according to a policy.

The key components are:

- subject (principal): the entity requesting access to resources

- object: a resource that can be accessed

- policy: rules that define how subjects can access objects

- authentication: identifying who is making the request

- authorization: determining what actions the subject is allowed to perform

- audit: logging actions performed by subjects

- accountability: holding subjects responsible for their actions

We formalize access decisions as:
\[
  \text{access}(\text{Subject}, \text{Object}) = \text{true} / \text{false}
\]

\subsection{Access Control Matrix}
A natural way to represent access control policies is through the access control matrix, which maps subjects to objects and specifies allowed actions.

\begin{table}[H]
  \centering
  \begin{tabular}{c|c|c|c|c|c}
      \toprule
       & \texttt{/user/alice} & \texttt{/user/bob} & \texttt{/etc/password} & \texttt{/share/bob/b.txt} & ... \\
    \midrule
      Alice & True & False & False & True &   \\
    \midrule
      Bob & False & True & False & True &   \\
    \midrule
      root & True & True & True & True &   \\
    \midrule
      ... &  &  &  &  &   \\
      \bottomrule
  \end{tabular}
  \caption{Access Control Matrix}
\end{table}

In practice, permissions are not simply binary. Instead, we use finer-grained access rights such as Read (R), Write (W), and Execute (X).

\begin{table}[H]
  \centering
  \begin{tabular}{c|c|c|c|c|c}
      \toprule
       & \texttt{/user/alice} & \texttt{/user/bob} & \texttt{/etc/password} & \texttt{/share/bob/b.txt} & ... \\
    \midrule
      Alice & \{R, W\} & \{\} & \{\} & \{R\} &   \\
    \midrule
      Bob & \{\} & \{R, W\} & \{\} & \{R, W\} &   \\
    \midrule
      root & \{R, W\} & \{R, W\} & \{R, W\} & \{R, W\} &   \\
    \midrule
      ... &  &  &  &  &   \\
      \bottomrule
  \end{tabular}
  \caption{Access Control Matrix}
\end{table}

However, directly storing this matrix is inefficient in real systems. This motivates practical mechanisms for implementing access control.

\section{Security Mechanisms}

There are two primary ways to implement the access control matrix:

1. Access Control List (ACL)

2. Capability-based Security

\subsection{Access Control List}
ACL is an object-centric mechanism. Each object maintains a list specifying which subjects have what permissions.

When access is requested:

- the system identifies the subject,

- looks up the subject in the object's ACL,

- and determines whether the requested action is allowed.

\subsection{Capability-based Security}
Capability-based security is a subject-centric mechanism. Each subject holds capabilities, which are unforgeable tokens that reference objects along with permitted operations.

When access is requested:

- the subject presents a capability,

- the system verifies the capability,

- and grants access if the required permission is included.

\subsection{Reference Monitor}
Regardless of the mechanism used, all access decisions must be enforced by a reference monitor.

A reference monitor mediates every access to objects and must satisfy:

1. unbypassable

2. tamper-proof

3. verifiable

Thus, subjects cannot access objects directly; all requests must pass through the reference monitor.

\subsection{Trusted Computing Base}
The Trusted Computing Base (TCB) consists of all components that must be correct to enforce security. Since all trust depends on the TCB, it should be:

- minimal

- unbypassable

- tamper-proof

- verifiable

\subsection{DAC vs. MAC}
Different systems adopt different policy models.

- Discretionary Access Control (DAC): object owners define access policies.

- Mandatory Access Control (MAC): a centralized authority enforces system-wide policies.

\subsection{Attribute-Based Access Control}
Attribute-Based Access Control (ABAC) generalizes access control by assigning attributes to subjects and objects. Policies are defined based on these attributes.

Role-Based Access Control (RBAC) can be viewed as a special case of ABAC.

\section{UNIX Security Model}

We now examine how these concepts are applied in a real system.

In UNIX:

1. subjects: users

2. objects: files, directories, devices, sockets, etc.

3. operations: read, write, execute

Users may belong to multiple groups, enabling group-based (role-like) access control.

\subsection{File Access Control}
Each file has an owner and a simple ACL based on:

1. owner

2. group

3. others

Permissions are set by the owner or root.

\begin{remark}
  If a user has write permission on a directory but not on a file, the ability to rename or delete the file depends on the sticky bit.

  - If the sticky bit is off: users can delete or rename files they do not own.

  - If the sticky bit is on: only the file owner, directory owner, or root can do so.
\end{remark}

Each process maintains:

1. real UID/GID: identity of the process owner

2. effective UID/GID: used for permission checks

3. saved UID/GID: allows privilege dropping and restoration

If an executable has the SUID bit set, executing it temporarily grants the process the privileges of the file owner.

\section{Privilege and Isolation}

Traditional UNIX distinguishes between:

- privileged processes (e.g., root)

- unprivileged processes

To reduce risk, modern systems decompose root privileges into smaller capabilities.

To limit damage from compromise:

- minimize the TCB

- enforce the principle of least privilege

- isolate untrusted components

\subsection{Sandbox}
A sandbox isolates a process from the rest of the system. If some vulnerable code is from an untrusted source or takes input from potential attackers, running it in a sandbox aims to prevent exploits from spreading outside of the sandbox. 

It restricts access to:

- file system

- network

- system resources

This prevents exploits from affecting the broader system.

\subsection{Android Process Isolation}
Android enforces sandboxing as follows:

1. Isolation: each app runs with a unique UID in its own VM

2. Interaction: permissions checked via a reference monitor

3. Least privilege: permissions granted explicitly by users

\subsection{Chrome Security Model}
Chrome uses a multi-process architecture:

- a privileged browser process

- multiple sandboxed renderer processes

Renderer processes handle untrusted web content with restricted permissions. If compromised, the damage is contained.

\subsection{Information Flow Tracking}
Information flow tracking monitors how data propagates through a system. Objects are labeled with security tags, which are updated during computation. This helps enforce confidentiality and integrity constraints.